\chapter{Résumé du Cours : L'Essentiel à Retenir}
\label{chap:resume}

Ce chapitre synthétise les concepts fondamentaux abordés dans les quatre modules du cours. Il sert de mémo pour les professionnels et de guide de révision pour les examens.

\section{Module 1 : Gouvernance et Analyse de Risque}
\textbf{Concept clé :} Le risque est le produit de la Probabilité par la Gravité. On ne peut pas éliminer tous les risques, on doit les gérer (Accepter, Réduire, Éviter, Transférer).
\begin{itemize}
    \item \textbf{ISO 27005 / EBIOS RM :} Méthodes pour identifier et traiter les risques.
    \item \textbf{RACI :} Outil indispensable pour définir qui fait, qui décide, qui consulte, qui informe.
    \item \textbf{RSSI :} Responsable ultime de la sécurité (Accountable), mais pas seul acteur.
\end{itemize}

\section{Module 2 : Audit et Management de la Qualité}
\textbf{Concept clé :} L'audit mesure l'écart entre ce qui est prévu (référentiel) et ce qui est fait (réalité). Il nécessite indépendance et preuves.
\begin{itemize}
    \item \textbf{ISO 19011 :} La norme "méthode" pour conduire un audit.
    \item \textbf{ISO 27001 :} La norme "exigence" pour le SMSI (certification).
    \item \textbf{Non-Conformité :} Écart constaté. Distinction Majeure (systémique) vs Mineure (isolée).
    \item \textbf{PAC :} Plan d'Actions Correctives (Correction immédiate + Action de fond).
\end{itemize}

\section{Module 3 : Gestion des Incidents et Continuité}
\textbf{Concept clé :} L'incident est inévitable. La résilience se mesure à la capacité de reprise (RTO/RPO).
\begin{itemize}
    \item \textbf{ISO 27035 :} Cycle de vie de l'incident (Détection, Confinement, Éradication, Rétablissement).
    \item \textbf{Forensics :} Investigation numérique. Règle d'or : ne rien altérer (Ordre de volatilité : RAM d'abord).
    \item \textbf{BIA (Business Impact Analysis) :} Analyse pour identifier les fonctions vitales.
    \item \textbf{RTO/RPO :} Objectifs de Temps de Reprise et Perte de Données acceptable.
    \item \textbf{PCA/PRA :} Plan de Continuité (Métier) vs Plan de Reprise (IT).
\end{itemize}

\section{Module 4 : Gestion de Crise et Retex}
\textbf{Concept clé :} La crise est un incident qui menace la survie de l'entreprise. Elle demande une gestion humaine et médiatique, pas seulement technique.
\begin{itemize}
    \item \textbf{Cellule de Crise :} Structure temporaire centralisée (Stratégique, Opérationnel, Communication).
    \item \textbf{Règle des 3 T :} Ténor (Porte-parole unique), Transparence, Temporalité (vitesse).
    \item \textbf{RETEX :} Retour d'expérience. Indispensable pour l'amélioration continue (PDCA). Distinction à chaud (immédiat) et à froid (analyse profonde).
\end{itemize}

% ==============================================================================
% Fichier : 07_evaluation_finale.tex
% Description : Evaluation Finale Globale
% ==============================================================================